\section*{Hoofdstuk \ref{ch:b1:structure-spectroscopy} --- Spectroscopie voor structuurbepaling}

\begin{solution}{exo:b1:structure-spectroscopy:1}
$\varepsilon = A/(\ell c) = 0.64/(1.00 \times \num{2.0e-5}) =
\qty{3.2e4}{L/(mol.cm)}$. Drie keer de concentratie: $A = 1.92$, als de
wet bij die \omterm{def:b1:structure-spectroscopy:absorbance}{absorbantie} nog geldt.
\end{solution}

\begin{solution}{exo:b1:structure-spectroscopy:2}
\ce{C6H6}: $D = (12 + 2 - 6)/2 = 4$, benzeen (een ring en drie
$\pi$-bindingen). \ce{C4H8O2}: 1, ethylethanoaat of butaanzuur.
\ce{C3H7NO}: $(6 + 2 + 1
- 7)/2 = 1$, propaanamide. \ce{C2H3Cl}: 1, chlooretheen. \ce{C10H14O}:
$(20 + 2 - 14)/2 = 4$: carvon heeft één ring, twee C=C en één C=O.
\end{solution}

\begin{solution}{exo:b1:structure-spectroscopy:3}
Propanon: één signaal (6H). Ethaanzuur: twee (3H, 1H). 2-Methylpropaan:
twee (9H, 1H). 1,2-Dichloorethaan: één (4H).
\end{solution}

\begin{solution}{exo:b1:structure-spectroscopy:4}
$\delta = 1460/400 = \qty{3.65}{ppm}$; bij \qty{90}{MHz} ligt het signaal
$3.65 \times 90 = \qty{329}{Hz}$ van de referentie. \omterm{def:b1:structure-spectroscopy:coupling}{Koppelingsconstanten}
blijven in hertz dezelfde, terwijl verschillen in verschuiving met het
veld groeien: \omterm{def:b1:structure-spectroscopy:coupling}{multipletten} die bij een laag veld overlappen, raken bij
een hoog veld gescheiden.
\end{solution}

\begin{solution}{exo:b1:structure-spectroscopy:5}
\ce{CH3}: triplet (3H); middelste \ce{CH2}: vijf buren, sextet
1\,:\,5\,:\,10\,:\,10\,:\,5\,:\,1 (2H); \ce{CH2Br}: triplet (2H), het
sterkst ontschermd.
\end{solution}

\begin{solution}{exo:b1:structure-spectroscopy:6}
Ethanol: O--H (3666) en geen C=O; propanon: C=O (1738) en geen O--H;
ethaanzuur: beide, O--H (3582) en C=O (1788).
\end{solution}

\begin{solution}{exo:b1:structure-spectroscopy:7}
Ethylethanoaat: kwartet (2H) rond 4.1, singlet (3H) rond 2.0, triplet
(3H) rond 1.3 ppm. Methylpropanoaat: singlet (3H) van \ce{OCH3},
ontschermd door de zuurstof (rond 3.7 ppm), kwartet (2H) van
\ce{CH2C=O} (rond 2.3), triplet (3H) rond 1.1. Het kwartet is in het
eerste het sterkst ontschermde signaal en in het tweede niet: een kwartet
rond 4 ppm betekent \ce{O-CH2CH3}.
\end{solution}

\begin{solution}{exo:b1:structure-spectroscopy:8}
Lijnen op \qty{0.080}{ppm} van elkaar, dat is $0.080 \times 90 = \qty{7.2}{Hz}$:
de \omterm{def:b1:structure-spectroscopy:coupling}{koppelingsconstante}. Het triplet heeft dezelfde afstand: de \ce{CH2}
en de \ce{CH3} zijn met elkaar gekoppeld.
\end{solution}

\begin{solution}{exo:b1:structure-spectroscopy:9}
Een C=O-band en één enkel signaal: propanon, \ce{CH3COCH3}. Het
aldehydische isomeer, propanal \ce{CH3CH2CHO}, vertoont de
C--H-strekvibratie van de \ce{CHO}-groep en zijn zeer sterk ontschermde
proton.
\end{solution}

\begin{solution}{exo:b1:structure-spectroscopy:10}
De eerste groep splitst de lijn in $n_1 + 1$ lijnen; elk daarvan wordt
door de tweede groep gesplitst in $n_2 + 1$: $(n_1 + 1)(n_2 + 1)$ lijnen.
Als $J_1 =
J_2$, hangt de positie van een lijn alleen af van $k_1 + k_2$, dat
$n_1 + n_2 +
1$ waarden aanneemt: de $n + 1$-regel met $n = n_1 + n_2$. Middelste
\ce{CH2} van 1-broompropaan: in het algemeen $4 \times 3 = 12$ lijnen,
zes (een sextet) wanneer de twee constanten gelijk zijn.
\end{solution}

\begin{solution}{exo:b1:structure-spectroscopy:11}
$D = 0$; een alcohol (O--H-band, geen C=O). Doublet (6H): twee
equivalente \ce{CH3} op één CH; \omterm{def:b1:structure-spectroscopy:coupling}{multiplet} (1H): die CH; doublet (2H): een
\ce{CH2} naast de CH die de OH draagt; brede singlet: OH.
2-Methylpropaan-1-ol, \ce{(CH3)2CH-CH2-OH}.
\end{solution}

\begin{solution}{exo:b1:structure-spectroscopy:12}
Elk deeltje absorbeert onafhankelijk: $A = A_1 + A_2 = \ell(\varepsilon_1c_1
+ \varepsilon_2c_2)$. Bij twee golflengten, met de vier coëfficiënten
bekend, geven twee lineaire vergelijkingen $c_1$ en $c_2$.
\end{solution}

\begin{solution}{pb:b1:structure-spectroscopy:1}
\textbf{1.} $0.5690 \times 12.0/44.0 = \qty{0.1552}{g}$.
\textbf{2.} $0.2328 \times 2.0/18.0 = \qty{0.02587}{g}$.
\textbf{3.} $0.2500 - 0.1552 - 0.0259 = \qty{0.0689}{g}$.
\textbf{4.} C \qty{62.1}{\percent}, H \qty{10.3}{\percent}, O
\qty{27.6}{\percent}.
\textbf{5.} C $0.1552/12.0 = 0.01293$, H $0.02587$, O $0.0689/16.0 =
0.00431$ mol: verhouding $3.00 : 6.00 : 1$, \ce{C3H6O}.
\textbf{6.} $M = \rho RT/p = 3740 \times 8.314 \times 373.15/(1.000 \times 10^5) =
\qty{116}{g/mol}$ (dichtheid in \unit{g/m^3}).
\textbf{7.} $116/58 = 2$: \ce{C6H12O2}, $D = (12 + 2 - 12)/2 = 1$.
\textbf{8.} Een C=O-strekvibratie.
\textbf{9.} Elke O--H: geen zuur, geen alcohol.
\textbf{10.} Een C--O-strekvibratie, sterk: de C--O van een ester.
\textbf{11.} C--H-strekvibraties van \ce{CH2}- en \ce{CH3}-groepen.
\textbf{12.} Eén $\pi$-binding, gebruikt door de C=O; met een sterke C--O
en geen O--H een ester. Een zuur heeft een O--H nodig; een diol heeft
geen C=O. Een keton met een ethergroep zou zijn C=O rond de
\qty{1738}{cm^{-1}} van propanon hebben, terwijl de waargenomen
\qty{1754}{cm^{-1}} dicht bij de \qty{1764}{cm^{-1}} van ethylethanoaat
ligt; de NMR van Deel III bevestigt de ester.
\textbf{13.} Vijf; $2 + 2 + 2 + 3 + 3 = 12$, de twaalf waterstofatomen.
\textbf{14.} Een \ce{CH2} met drie buren (een \ce{CH3}), gebonden aan de
zuurstof (ontschermd): \ce{-O-CH2-CH3}.
\textbf{15.} De \ce{CH3} van die ethylgroep, met twee buren; de twee
signalen zijn met elkaar gekoppeld, vandaar één $J$.
\textbf{16.} Een \ce{CH2} met twee buren, naast de C=O.
\textbf{17.} Een \ce{CH2} met vijf buren: tussen een \ce{CH2} en een
\ce{CH3}.
\textbf{18.} Een \ce{CH3} met twee buren, ver van de karakteristieke
groep.
\textbf{19.} \ce{CH3CH2O-} en \ce{-CH2CH2CH3}.
\textbf{20.} De protonen bij 4.13 ppm zitten op het koolstofatoom dat aan
zuurstof gebonden is, die bij 2.28 ppm naast de C=O: beide groepen zijn
elektronenzuigend.
\textbf{21.} \ce{CH3CH2-O-CO-CH2CH2CH3} of \ce{CH3CH2-CO-O-CH2CH2CH3}.
Alleen de eerste plaatst de \ce{CH2} met drie buren op de zuurstof (het
kwartet bij 4.13 ppm).
\textbf{22.} Propylpropanoaat is de tweede: zijn \ce{OCH2} zou een triplet
zijn en het kwartet zou rond 2.3 ppm liggen.
\textbf{23.} Butylethanoaat, \ce{CH3CO-O-C4H9}, zou een singlet (3H) rond
2.0 ppm vertonen en geen kwartet.
\textbf{24.} Ethylbutanoaat, \ce{CH3CH2CH2-CO-O-CH2CH3}.
\textbf{25.} IR: C=O 1754, C--O 1182, C--H 2982, geen O--H. NMR: \ce{OCH2}
4.13 (q), \ce{CH2CO} 2.28 (t), middelste \ce{CH2} 1.66 (sextet),
\ce{CH3} van de ester 1.25 (t), \ce{CH3} van de keten 0.95 (t),
\omterm{def:b1:structure-spectroscopy:integration}{integraties} 2, 2, 2, 3, 3.
\textbf{26.} \textbf{\qty{116}{g/mol}}: ethylbutanoaat, \ce{C6H12O2}.
\end{solution}
